\section{Evidence and reproducible replay}
\label{sec:replay}

\subsection{What the saved compiler record establishes}
The inherited run is \code{verification/final_finite_release_run/SUMMARY.json} in the unchanged v0.6.0 snapshot. It began on 6 September 2026 at 15:30:24 UTC and was updated at 15:32:49.641452 UTC. It records successful compilation of all 25 current modules, all 323 named theorem axiom audits, and two deliberate rejected controls. Current static reconciliation checks the source hashes and saved log hashes against those records.

The allowed axiom set is \code{propext}, \code{Classical.choice}, and \code{Quot.sound}. Every reported theorem dependency set is a subset of it. One reported theorem has no axiom dependencies. These are foundational dependencies accepted in the recorded Lean environment. They do not remove manuscript hypotheses such as invertibility of the full resolvent or Kraus normalization. Axiom auditing and hypothesis reading answer different questions.

The separate inherited \code{CONSTRUCTED_OBJECTS_AUDIT.json}, in the directory for the same release run, records six constructed objects: the fiber, projection to the fiber, local product homeomorphism, local continuous linear fiber equivalence, covariant derivative taking values in the fiber, and fiber curvature. Definitions can contain proofs of inverse, linearity or continuity fields, so an inventory of theorem names alone would miss part of this evidence. The audit of six objects remains separate from the count of 323 theorems. The computation report dated 9 October 2026 records 12 correspondence checks among the inherited audit, source, inventory, runner, transcript and cumulative summary. Those checks added no theorem declarations and did not execute the Lean compiler.

The source archive omits the compiled public and private interfaces, the portable compiler kit and the Mathlib cache. The saved success record is inherited compilation evidence. It is not a fresh compilation on this host or an independent recheck of all imported Mathlib proofs. The recorded release reused its pinned library environment and did not rebuild earlier papers.

\subsection{Source and document checks in this edition}
The guide resolves nested namespaces in active source and derives 323 distinct theorem declarations across 25 modules. That set matches the saved theorem inventory and every saved theorem axiom report. All 19 finite ledger entries resolve, including \code{GlobalConnection.covariant} among the geometric supporting endpoints. Where a ledger entry gives only a module, the guide adds explicitly identified source pointers. It preserves the original ledger.

All 67 selected excerpts are exact contiguous source blocks and retain exactly the bytes of the retrieved guide's raw UTF-8 selections. Their line ranges, original module hashes and excerpt hashes accompany this document. Rendered code uses explicit font fallback for the mathematical script letters and small square symbol used for dot products absent from the monospaced font. The raw source is unchanged. The PDF checks cover every theorem name, excerpt, reference, missing glyph and text bound.

The local \code{static_check.py} checks saved source and log identities, transitive local import hashes, declarations, dependency pins, inherited rejection records and the audit of six constructed objects. It also compares every delimited mathematical span in the six teaching chapters with the retrieved edition. The prose is checked against a preserved baseline and an explicit record of the approved changes. That record covers each complete chapter, including its hypotheses and explanations of proofs. These checks establish editorial preservation and consistency with the source. They do not compile Lean or validate a new proof.

\subsection{The finite computations recorded on 9 October 2026}
The report \code{verification/monograph_checks/COMPUTATIONAL_REVIEW.json} records the finite computations performed on 9 October 2026. Its six records retain the commands, environment versions, limits, inputs, outputs and hashes of those executions. The original and corrected numerical checks used Python 3.12.14, NumPy 2.3.5 and SciPy 1.17.0. The symbolic checks used SymPy 1.14.0. Statements about unavailable dependencies in the retained reports of 6 October describe the host used at that earlier date. The present proofreading pass preserves the October 9 evidence without repeating those numerical or symbolic executions.

In that recorded run, the original companion passed 22 pytest cases. A second entry point called the same 22 functions directly, so the report contains 22 distinct tests. The verifier software passed 16 regression methods, comprising nine positive methods and seven controls using malformed inputs or tampering. These Python controls differ from the two inherited deliberate Lean rejections. No fresh kernel run or compiler rejection test was possible during this proofreading pass because \code{lean}, \code{lake} and \code{elan} are absent.

The report records an exact match for all 1,048 stored result leaves, including 1,045 numerical fields, against the original experiment replay. The recorded comparison bounds are absolute $10^{-10}$ and relative $10^{-8}$, although equality made those allowances unnecessary in the recorded run. The unchanged strengthened suite passed 18,049 checks in 15 groups. Its separately corrected copy passed 17,699 checks in the same 15 groups, with 87 exact mode-index assertions for $N=4,\ldots,32$. The two lower endpoints were changed from \code{-N//2} to \code{-(N//2)}. For odd \code{N}, Python floor division otherwise adds a repeated residue representative. The difference of 350 checks reflects corrected bookkeeping. The raw source remains intact, and the two overlapping suite totals are not added.

The independent exact rational witnesses contain 269 assertions in total, including three explicit negative controls and a separate singular-inverse rejection benchmark. The original checks of visual examples executed 180 SymPy assertions, including nine rejected controls and 171 positive assertions. The extended checks executed 373 positive exact assertions and 12 rejected controls over finite examples using rational and Gaussian rational numbers. All 38 vector plate sources regenerated with identical bytes, comprising 26 original plates and 12 extended plates. Regenerating the plates checks consistency of the presentation and does not add mathematical proof coverage.

A separate audit, \code{verification/monograph_checks/TRADITIONAL_EXACT_MODELS_REVIEW.json}, records checks of the three traditional illustrations added in that edition. It passed 30 positive exact assertions and rejected six deliberately false assertions. The model JSON and three generated PDFs reproduced with identical bytes. These are fixed examples of seed support, block inversion and channel scaling. They add no formal targets or general recovery conclusion.

The manifest records retain the actual finite ranges, conventions and tolerances. Original numerical checks typically use bounds from $10^{-12}$ to $10^{-10}$, with a $10^{-6}$ certificate for moving frames. The strengthened base tolerance is $2\times10^{-9}$, with individual bounds through $2\times10^{-7}$. Exact arithmetic gates have no numerical tolerance. These executions overlap and use different count units. Their totals are neither summed as independent evidence nor interpreted as universal theorems. They do not settle any of the five physical continuum targets.

\begin{longtable}{@{}p{.31\linewidth}p{.30\linewidth}p{.32\linewidth}@{}}
\toprule Evidence & Time or source & Meaning\\\midrule\endhead
25 module compiler exits & Saved 6 September release & Inherited compilation under the locked toolchain.\\
323 theorem axiom records & Saved release, currently reconciled & Current source declarations agree with the stored audited set.\\
Audits of six constructed objects & Separate saved release record & Proof fields in definitions are audited separately.\\
Source, excerpt and document checks & Proofreading on 18 October & Exact source correspondence, reviewed prose and readable navigation.\\
22 pytest cases, 16 verifier methods & Recorded 9 October execution & Software checks at the recorded inputs and conditions.\\
Numerical and exact witness checks & Recorded 9 October executions & Finite computations with explicit arithmetic and assumptions.\\
38 vector source regenerations & Recorded 9 October replay & Deterministic plate sources, with no added theorem coverage.\\
\bottomrule\end{longtable}

\subsection{Rebuild only this guide}
The guide directory contains the authored TeX, frozen navigation, rendered code, raw excerpts, styles and font licenses. Font binaries are not included in this delivery. The styles use the matching bundled fonts when those are supplied and otherwise use the named installed fonts. The script \code{build.py} copies the complete source tree to an isolated temporary directory. It performs four XeLaTeX passes to resolve the document structure, then runs \code{xdvipdfmx} once to produce the PDF. It compiles the authored document without regenerating prose, source data, figures or mathematical evidence, and it does not invoke Lean.

\begin{Verbatim}[fontsize=\small,breaklines=true,breakanywhere=true]
python editable/lean_guide/build.py
python editable/lean_guide/build.py --output-directory /path/to/pdf_rebuild
python editable/lean_guide/build.py --check-only
\end{Verbatim}

Run these commands from the package root, or supply an absolute path to the script. Python 3, XeLaTeX, \code{xdvipdfmx} and the TeX packages named in the preamble are required. Fonts resolve through the relative \code{fonts/} directory when the matching binaries are present. Otherwise the styles use installed P052, Nimbus Sans, Latin Modern and DejaVu Sans Mono fonts. The optional \code{--output-directory} keeps delivered products unchanged, and \code{--check-only} writes no delivered products.

Regenerating the source tables additionally requires FontTools and the preserved snapshot. The assembler uses the bundled DejaVu Sans Mono and Latin Modern Math files when present, or matching installed files otherwise. Run the following commands only when regenerating the tables as part of an edited release. Refresh the recorded checks and document hashes afterward.

\begin{Verbatim}[fontsize=\small,breaklines=true,breakanywhere=true]
python editable/lean_guide/assemble.py
python editable/lean_guide/static_check.py
python editable/lean_guide/build.py
python editable/lean_guide/verify_pdf.py
\end{Verbatim}

Both source scripts accept \code{--source PATH} for a relocated snapshot. Assembly stops if counts, inventory, saved axiom sets or source identities disagree. A new source edition requires a deliberate revision of the expected counts and endpoint selection together with its own evidence. These scripts reconcile documentation. They are not proof checkers.

\subsection{Replay the formal development with the pinned kit}
With the separately supplied pinned portable kit, set the working directory to \code{source/complete/Paper_VII_Complete_v0.6.0_With_Visual_Appendix/original_snapshot}, relative to the package root. Choose a new output directory to preserve the inherited evidence.

\begin{Verbatim}[fontsize=\small,breaklines=true,breakanywhere=true]
python scripts/verify.py --kit /path/to/kit \
  --out verification/my_run --timeout 240
python scripts/audit_constructions.py --kit /path/to/kit \
  --run verification/my_run
python scripts/test_verifier.py
\end{Verbatim}

The verifier checks runtime version and Mathlib revision, then compiles local modules in dependency order with one worker and strict warnings. It compares emitted axiom audits with the exact inventory and checks both deliberate invalid controls. The construction audit uses the fresh run's actual build interfaces. It rejects stale source or interface identities.

For short execution windows, the retained checkpoint mechanism can pause after a few checks:

\begin{Verbatim}[fontsize=\small,breaklines=true,breakanywhere=true]
python scripts/verify.py --kit /path/to/kit \
  --out verification/my_run --timeout 240 \
  --max-checks 3 --resume
\end{Verbatim}

Repeat until the verifier reports complete success. Exit code 2 means paused or incomplete. Resume checks the verifier, runner, inventory, source dependencies and both public and private interfaces before trusting a cached module. An interrupted build remains incomplete even when earlier modules passed.

The \code{check_low_memory.py} launcher uses ordinary Lean with one worker. Its name does not promise an enforced memory cap. The snapshot is not a delivered Lake project with a current \code{lakefile}. Its provenance manifest supplies the portable kit's locked dependencies. A bare \code{lake build} would not reproduce the retained procedure.

\subsection{The wrapper in the retained complete bundle}
Set the working directory to \code{source/complete/Paper_VII_Complete_v0.6.0_With_Visual_Appendix}. Its wrapper provides three actions:

\begin{Verbatim}[fontsize=\small,breaklines=true,breakanywhere=true]
python Run-PaperVII.py verify
python Run-PaperVII.py examples --output-dir /path/to/example_replay
python Run-PaperVII.py original --output-dir /path/to/original_replay
\end{Verbatim}

The integrity action verifies the manifest and checksums. The example action runs the original example suite for the visual appendix with its declared dependencies. The original action runs the original workflow for numerical examples, figures, TeX and verifier software. Its own documentation identifies that action as a companion replay rather than Lean execution. The October 9 report records the individual checks that were actually executed. It does not claim that the original combined PowerShell procedure, the full figure generation procedure or the historical manuscript rebuild ran as a single workflow.

All nine dependency revisions and the compiler commit appear in the dependency appendix. Source and transcript hashes bind this reading to its inherited artifacts. Document construction, finite computation and kernel replay each have a distinct purpose. The release records their outcomes separately.
